\par\noindent\SourceLeaf{55}%
\hspace*{\parindent}Thus a very complicated law appears to admit a very short expression in terms of binomial coefficients.
For although some factors of one cancel against factors of the other when these coefficients are expanded into products, as they can be, the resulting expression is much longer than the earlier one. To verify this in practice, for the arbitrary $r^{\text{ti}}$ coefficient we give the general term $h^{\text{tus}}$ in the shortest form this method affords

\[
\frac{\nu.(\nu-1).\,.\,.(\nu-(r+2h-1))}{1.2^2.3^2.\,.\,. h^2} . \frac{1}{(h+1)(h+2).\,.\,.(h+r)}\frac{1}{2^{r+2h-1}}n^{r+2h}.
\]

\textls[120]{Note.}
Anyone wishing to deduce the general law from the series that \textsc{Euler}\HistoricalNote{*)}{*)}{\textsc{Euler}, loc. cit., Problem 36.} proposed for some of the initial coefficients will obtain a formula of this kind. The approach that the illustrious author indicates for investigating that law seems less direct and less easy for obtaining the series already set out. In another respect, however, it seems excellent and well worth considering.
The whole procedure consists in setting up for the quantity $(1+n\cos\varphi)^{\nu}$ the series

\[
A+A^{(1)}\cos\varphi+A^{(2)}\cos2\varphi+.\,.\,.\,.
\]

\noindent and then using differentiation to derive the recurrence scheme that governs the dependence of its coefficients on one another. Since it involves only three terms, it gives any coefficient from the two preceding it. Thus, provided that the first two are given in finite form, all the others can be expressed in finite form. Although this relationship among the coefficients can be derived from the series already found, it is more easily established by \textsc{Euler}'s method.

\Needspace{8\baselineskip}
7. Let it be, as before,

\[
(1+n\cos\varphi)^{\nu}=A+A^{(1)}\cos\varphi+A^{(2)}\cos2\varphi+.\,.\,.\,.+A^{(h)}\cos h\varphi+.\,.\,.
\]

\noindent Taking the differentials of the logarithms on both sides and cancelling $d\varphi$ gives

\[
\frac{\nu n\sin\varphi}{1+n\cos\varphi}=\frac{A^{(1)}\sin\varphi+2A^{(2)}\sin2\varphi+3A^{(3)}\sin3\varphi+.\,.\,.+hA^{(h)}\sin h\varphi+.\,.\,.}{A+A^{(1)}\cos\varphi+A^{(2)}\cos2\varphi+.\,.\,.+A^{(h)}\cos h\varphi+.\,.\,.}
\]

$\alpha$) Clearing the denominators by multiplication gives on one side

\[
\begin{aligned}
&\nu n\sin\varphi\bigl(A+A^{(1)}\cos\varphi+A^{(2)}\cos2\varphi+.\,.\,.\,.+A^{(h)}\cos h\varphi\\&\quad+A^{(h+1)}\cos(h+1)\varphi+A^{(h+2)}\cos(h+2)\varphi+.\,.\,.\,.\bigr).
\end{aligned}
\]
